% Copyright (c) 2026 Eugene Joseph M. Ragasa.
% SPDX-License-Identifier: Apache-2.0
% Author: Eugene Joseph M. Ragasa
% Formal proof completed: 2026-10-19
\documentclass[11pt]{article}
\usepackage[T1]{fontenc}
\usepackage{amsmath,amssymb,booktabs}
\usepackage[margin=1in]{geometry}
\usepackage[hidelinks]{hyperref}
\usepackage{tabularx}
\title{Proof Dossier: Hilbert-Ball Intersection Criterion}
\author{Eugene Joseph M. Ragasa}
\date{Formal proof completed: 2026-10-19}
\begin{document}
\maketitle
\section{Identity and status}
\begin{tabularx}{\textwidth}{@{}l>{\raggedright\arraybackslash}X@{}}
\toprule
Statement ID & \texttt{BG-HILBERT-THM-001} \\
Status & \texttt{proved} \\
Lean declaration & \path{ProjectKoios.BandGauge.Hilbert.hilbertAdmissibleBalls_inter_nonempty_iff} \\
Owning module & \path{ProjectKoios/BandGauge/Hilbert/AdmissibleBalls.lean} \\
\bottomrule
\end{tabularx}

\section{Exact statement}
For centers $c_1,c_2$ in a real Hilbert space and radii $r_1,r_2\ge0$,
\[
 B(c_1,r_1)\cap B(c_2,r_2)\ne\varnothing
 \quad\Longleftrightarrow\quad
 d(c_1,c_2)\le r_1+r_2.
\]

\section{Source and scope}
This is classical normed-space geometry. The motivating owner source is
\path{ksdft2effmass} at immutable commit
\begin{center}
\small\texttt{99a46cb59d0c6f27ada7d709f194996ef9619d9f}.
\end{center}
The finite weighted Frobenius loss appears in the Paper A v3 manuscript, lines
575--601. Its repository path is
\path{docs/publications/papers/Ragasa2026a/v3/}\path{manuscript.tex}. The relationship
is \texttt{background-only}; this theorem does not formalize that loss.

\section{Proof}
A common point gives the forward inequality by the triangle inequality. Conversely, if
$c_2$ already lies in the first ball, use $c_2$. Otherwise choose the point at fraction
$r_1/d(c_1,c_2)$ along the line from $c_1$ to $c_2$. Its distances from the centers are
$r_1$ and $d(c_1,c_2)-r_1\le r_2$.

\section{Boundary and trusted base}
No matrix representation, Frobenius weights, Bloch boundary condition, operator norm,
or numerical claim is present. The Lean axiom report contains \texttt{propext},
\texttt{Classical.choice}, and \texttt{Quot.sound}. No admission or scientific axiom is
used. The proof and prose were AI-assisted and kernel-checked. No independent expert
review is recorded. Status: \texttt{proved}; working draft.
\end{document}
